\section{Studi Pendahuluan}

Studi pendahuluan diawali dengan analisis karakteristik dataset yang dilakukan sebelum eksperimen apa pun dijalankan.
Analisis ini memeriksa properti tiap dataset yang menentukan kecocokannya dengan tiap hipotesis, seperti skala dan dimensi ruang solusi D, jumlah \textit{permit} unik sebagai indikator risiko duplikasi, rasio \textit{permit}, serta lantai derau divergensi Jensen-Shannon antar-jendela pada kondisi stasioner sebagai penentu kelayakan mekanisme pemicu-\textit{drift}, sehingga keputusan penggunaan tiap dataset didasarkan pada bukti karakteristik dan bukan asumsi.
Kelima dataset ABAC Lab beserta dataset validasi eksternal \textit{Smart Building IoT} diprofilkan dengan prosedur yang sama yang hasilnya dirangkum pada Tabel \ref{tab:karakteristik-dataset}.

\begin{longtable}{|>{\raggedright\arraybackslash}p{2cm}|>{\raggedright\arraybackslash}p{1.5cm}|>{\raggedright\arraybackslash}p{2.5cm}|>{\raggedright\arraybackslash}p{1.5cm}|>{\raggedright\arraybackslash}p{2.5cm}|>{\raggedright\arraybackslash}p{3cm}|}
\caption{Karakteristik Dataset dan Kecocokannya dengan Hipotesis (Analisis Pra-Eksperimen)}
\label{tab:karakteristik-dataset} \\
\hline
\textbf{Dataset} & \textbf{Dimensi $D$} & \textbf{Permit unik (risiko duplikasi)} & \textbf{Rasio permit} & \textbf{Lantai derau JSD\label{first:jsd} (H3)} & \textbf{Kecocokan hipotesis} \\
\hline
\endfirsthead

\multicolumn{6}{c}{{\bfseries \tablename\ \thetable{} -- lanjutan dari halaman sebelumnya}} \\
\hline
\textbf{Dataset} & \textbf{Dimensi $D$} & \textbf{Permit unik (risiko duplikasi)} & \textbf{Rasio permit} & \textbf{Lantai derau JSD (H3)} & \textbf{Kecocokan hipotesis} \\
\hline
\endhead

\hline
\multicolumn{6}{r}{Lanjut ke halaman berikutnya} \\
\endfoot

\hline
\endlastfoot

% ==================== DATA BARIS ====================
workforce & 436 & 14.513 (rendah) & 0{,}30 & 0{,}0066 (layak) & H1 baik; H2/H4 baik \\
\hline
e-document & 1.907 & 32.961 (rendah) & 0{,}30 & 0{,}0325 (tak layak) & H1 baik; H3 terbatas (lantai derau $>$ $\tau$); H2/H4 baik \\
\hline
healthcare & 46 & 14 (TINGGI) & 0{,}30 & 0{,}0022 (layak) & H1 tumpul (duplikasi); H2/H3/H4 baik \\
\hline
project-management & 51 & 17 (TINGGI) & 0{,}30 & 0{,}0010 (layak) & H1 tumpul (duplikasi); H2/H3/H4 baik \\
\hline
university & 54 & 90 (TINGGI) & 0{,}30 & 0{,}0026 (layak) & H1 tumpul (duplikasi); H2/H3/H4 baik \\
\hline
Smart Building IoT* & 428* & 5.729 (rendah) & 0{,}75 & 0{,}0065 (layak) & berat-permit (sepele F1=0{,}857); H1 vs garis sepele; H2/H3/H4 baik \\
\hline

\end{longtable}

Tanda (*) menunjukkan dimensi efektif \textit{Smart Building IoT} setelah subsampel terstratifikasi 8.000 baris dan penyatuan nilai ber-\textit{support} rendah (berkas asalh 5,7 juta permintaan, \textit{permit} 4.273.814 dan \textit{deny} 1.429.088).
Rasio \textit{permit} seragam 0,30 pada dataset ABAC Lab merupakan parameter konstruksi log, sedangkan 0,75 pada \textit{Smart Building IoT} adalah rasio bawaan datanya.

Pembacaan Tabel \ref{tab:karakteristik-dataset} menghasilkan tiga implikasi rancangan yang memandu penggunaan dataset pada eksperimen.
Pertama, untuk H1, hanya \textit{workforce} dan \textit{e-document} yang memiliki \textit{permit} unik memadai sehingga pembeda antar-algoritma tidak tumpul.
Kedua, untuk H2, keenam dataset sama-sama layak karena kedataran memori terhadap ukuran populasi merupakan properti algoritma yang tidak bergantung pada karakteristik dataset.
Ketiga, untuk H3, kelayakan justru terbalik dengan H1 karena \textit{e-document} tidak layak karena lantai derau JSD-nya (0,0325) melampaui ambang $\tau$ = 0,01 sehingga pemicu \textit{drift} menyala di setiap jendela, sedangkan ketiga dataset kecil maupun \textit{Smart Building IoT} berlantai derau rendah (0,0010-0,0065 < $\tau$) sehingga pemicu dapat diam saat data stabil.
Perlu dicatat bahwa nilai absolut lantai derau ini bukan konstanta melainkan bergantung pada ukuran jendela pengukuran, karena bias sampel-hingga estimator informasi berbanding lurus dengan rasio jumlah sel kontingensi terhadap jumlah sampel per jendela.
Pada jendela analisis ini (800 entri), lantai derau \textit{e-document} terukur 0,0325, sedangkan pada jendela eksperimen \textit{drift} yang lebih besar (2.000 entri) ia turun ke kisaran 0,013-0,016.
Keduanya tetap di atas $\tau$ = 0,01, sehingga vonis ketidaklayakan H3 pada \textit{e-document} kokoh terhadap pilihan ukuran jendela.
Implikasi pentingnya adalah dataset \textit{Smart Building IoT} menjadi sarana untuk menunjukkan mekanisme H3 yang tidak dapat diperlihatkan pada \textit{e-document}, sehingga penggunaan seluruh dataset untuk seluruh hipotesis (bukan hanya dua dataset utama) memperoleh justifikasi empiris jadi setiap dataset menyumbang pada hipotesis yang paling sesuai dengan karakteristiknya.

Analisis ini dijalankan dengan instrumen tersendiri di luar rantai algoritma yang diuji, sehingga tidak memengaruhi anggaran evaluasi maupun kesetaraan algoritma.
Instrumen ini berfungsi murni sebagai penyaring kecocokan pra-eksperimen.
Setelah karakteristik dataset dipahami, rantai metodologi diuji melalui studi pendahuluan algoritmik yang dilaporkan pada paragraf-paragraf berikut.

Sebelum eksperimen penuh dilaksanakan, seluruh rantai metodologi telah diuji melalui studi pendahuluan pada perangkat target.
Studi ini semula dimaksudkan untuk memverifikasi kebenaran implementasi dan mengalibrasi parameter eksekusi, tetapi menghasilkan temuan yang berimplikasi langsung pada temuan hipotesis.
Subbab ini melaporkan palksanaan dan hasilnya secara lengkap, termasuk hasil yang tidak mendukung rumusan awal, agar dasar revisi hipotesis dapat ditelusuri dan dinilai secara terbuka.
Dengan demikian studi pendahuluan ini diperlakukan sebagai tahap eksploratori yang menyempurnakan rumusan H1, H3, dan H4.
Untuk mencegah hipotesis dibentuk dan diuji oleh data yang sama, eksperimen final pada Bab IV memakai kumpulan \textit{seed} (baik \textit{seed} pembangkitan dataset maupun \textit{seed} algoritma) yang berbeda dari \textit{seed} yang dipakai pada studi pendahuluan ini sehingga penyempurnaan rumusan pada tahap eksplorasi tidak mengambil manfaat dari data yang kemudian dipakai untuk mengujinya.

Studi pendahuluan membandingkan CoDA-EDA terhadap dua \textit{baseline} keluarga \textit{compact} (\textit{compact-GA} dan \textit{compact-PSO}) pada kedua dataset utama, mengikuti seluruh kontrol kesetaran, yakni representasi biner, fungsi \textit{fitness}, operator perbaikan, prosedur pembenihan, dan anggaran evaluasi total yang identik.
Empat dataset ber-\textit{seed} berbeda dijalankan dengan sepuluh replikasi algoritma pada masing-masingnya dan menghasilkan 40 replikasi per kondisi.
Anggaran evaluasi diskalakan menurut dimensi soal yakni 12.001 evaluasi untuk \textit{workforce} dan 52.001 untuk \textit{e-document}.
Pengujian mengikuti protokol yang ada, yaitu uji Friedman sebagai uji omnibus, dilanjutkan uji Wilcoxon signed-rank berpasangan dengan koreksi Holm-Bonferroni bila uji omnibus signifikan, diserta pelaporan \textit{effect size rank biserial}.

\begin{longtable}{|p{2cm}|c|p{1.5cm}|p{1.5cm}|p{1.5cm}|>{\raggedright\arraybackslash}p{2cm}|>{\raggedright\arraybackslash}p{3cm}|}
\caption{Hasil Studi Pendahuluan: Kualitas Kebijakan (F-score) pada Dua Dataset Utama}
\label{tab:studi-pendahuluan-fscore} \\
\hline
\textbf{Dataset} & \textbf{Korelasi} & \textbf{CoDA-EDA} & \textbf{compact-GA} & \textbf{compact-PSO} & \textbf{Friedman} & \textbf{Post-hoc CoDA-EDA vs compact-GA} \\
\hline
\endfirsthead

\multicolumn{7}{c}{{\bfseries \tablename\ \thetable{} -- lanjutan dari halaman sebelumnya}} \\
\hline
\textbf{Dataset} & \textbf{Korelasi} & \textbf{CoDA-EDA} & \textbf{compact-GA} & \textbf{compact-PSO} & \textbf{Friedman} & \textbf{Post-hoc CoDA-EDA vs compact-GA} \\
\hline
\endhead

\hline
\multicolumn{7}{r}{Lanjut ke halaman berikutnya} \\
\endfoot

\hline
\endlastfoot

% ==================== DATA BARIS ====================
Workforce ($D=436$) & Rendah & 0{,}8354 & 0{,}8409 & 0{,}8123 & $p = 0{,}025$ & tidak signifikan ($p = 0{,}545$) \\
\hline
Workforce ($D=436$) & Sedang & 0{,}8644 & 0{,}8486 & 0{,}8471 & $p = 0{,}592$ & tidak dilanjutkan (omnibus tidak signifikan) \\
\hline
Workforce ($D=436$) & Tinggi & 0{,}9017 & 0{,}8852 & 0{,}8769 & $p = 0{,}161$ & tidak dilanjutkan (omnibus tidak signifikan) \\
\hline
E-document ($D=1.907$) & Rendah & 0{,}3754 & 0{,}2071 & 0{,}2395 & $p = 8{,}5 \times 10^{-7}$ & signifikan; $p_{\text{adj}} = 1{,}1 \times 10^{-4}$; $r = +0{,}751$ \\
\hline
E-document ($D=1.907$) & Sedang & 0{,}3534 & 0{,}2308 & 0{,}2685 & $p = 1{,}9 \times 10^{-4}$ & signifikan; $p_{\text{adj}} = 4{,}9 \times 10^{-4}$; $r = +0{,}685$ \\
\hline
E-document ($D=1.907$) & Tinggi & 0{,}4160 & 0{,}2944 & 0{,}3067 & $p = 1{,}8 \times 10^{-3}$ & signifikan; $p_{\text{adj}} = 7{,}9 \times 10^{-4}$; $r = +0{,}663$ \\
\hline

\end{longtable}

Hasil pada Tabel \ref{tab:studi-pendahuluan-fscore} memperlihatkan dua pola yang bertentangan dengan rumusan awal.
Pertama, pada \textit{workforce} keunggulan CoDA-EDA atas \textit{compact-GA} hanya sebesar +0,016 pada korelasi sedang maupun tinggi dan tidak mencapai taraf signifikansi (p = 0,197 dan p = 0,111 pada uji langsung).
Uji tren pertumbuhan keunggulan seiring naiknya korelasi juga tidak signifikan (p = 0,073).
Kedua, pada \textit{e-document} keunggulan CoDA-EDA justru sangat besar dan signifikan di seluruh tingkat korelasi dengan selisih \textit{F-score} +0,12 hingga +0,17 dan \textit{effect size} besar (r = 0,66-0,75), namun keunggulan tersebut tidak tumbuh seiring korelasi dan selisih terbesar justru teramati pada korelasi rendah (p = 0,506 untuk uji tren).
Dengan demikian rumusan awal H1 yang memprediksi keunggulan muncul khusus pada korelasi tinggi dan mendekati \textit{baseline} pada korelasi rendah tidak didukung data pada kedua dataset.

Faktor yang secara konsisten membedakan kedua dataset adalah dimensi ruang solusi.
Pada \textit{workforce} (D = 436), \textit{compact-GA} yang murni univariat sudah mencapai \textit{F-score} 0,885 pada korelasi tinggi sehingga ruang perbaikan yang tersisa sangat tipis.
Pada \textit{e-document} (D = 1.907), kualitas \textit{compact-GA} turun ke 0,294 sementara CoDA-EDA bertahan pada 0,416.
Penjelasan mekanistiknya sejalan dengan teori EDA, yakni ketika ruang pencarian membesar, \textit{sampling} yang diarahkan struktur dependensi menjadi jauh lebih efektif dibanding \textit{sampling} univariat yang harus menjelajahi ruang berdimensi tinggi tanpa panduan keterkaitan antar-variabel.
Atas dasar inilah H1 dirumuskan ulang dengan dimensi ruang solusi sebagai faktor penentu dan dimensi D ditambahkan sebagai variabel bebas.

Studi pendahuluan juga mengonfirmasi dua hal yang memperkuat rancangan.
Pertama, kemampuan memanen korelasi memang berbeda antar algoritma.
Kenaikan \textit{F-score} dari korelasi rendah ke tinggi pada \textit{workforce} mencapai $+0{,}066$ bagi CoDA-EDA ($p = 1{,}6 \times 10^{-6}$; $r = +0{,}873$) dibanding $+0{,}044$ bagi \textit{compact-GA} ($p = 1{,}0 \times 10^{-3}$; $r = +0{,}598$), sehingga premis mekanistik mengenai manfaat pemodelan dependensi tetap terdukung meskipun bentuk H1 semula tidak.
Kedua, CoDA-EDA secara konsisten menghasilkan kebijakan yang lebih ringkas.
WSC rata-rata 23,8 berbanding 28,9 (\textit{compact-GA}) dan 38,2 (\textit{compact-PSO}) pada \textit{workforce}.
Keringkasan ini merupakan dimensi kualitas kebijakan yang sah dan relevan bagi administrator, namun belum tercakup pada rumusan hipotesis awal.

Rangkaian tahap studi pendahuluan berikutnya (perluasan ke tujuh algoritma, pengukuran jejak memori, simulasi siklus pembaruan lintas jendela, pengujian pada \textit{e-document}, eksplorasi perbaikan cakupan, dan validasi eksternal) berimplikasi langsung pada perumusan hipotesis.
Ringkasnya H1 mengalami tiga kali revisi (dari bergantung tingkat korelasi, ke bergantung dimensi ruang solusi, hingga klaim terbatas terhadap keluarga \textit{compact} univariat pada anggaran aturan memadai); H3 diberi syarat tingkat \textit{drift}; dan H4 diperlonggar menjadi sekurang-kurangnya setara.

Tahap-tahap tersebut sekaligus menghasilkan kalibrasi rancangan dan pengukuran yang menjadi ketetapan bagi eksperimen.
Penetapan ambang deteksi \textit{drift} $\tau$ = 0,01 dari pemisahan sinyal divergensi Jensen-Shannon terhadap lantai derau yang membesar seiring jumlah sel kontingensi; kalibrasi bobot fungsi \textit{fitness} $w_\text{fp}$\label{first:symbol-wfp} melalui uji monotonitas; serta optimasi implementasi perhitungan \textit{mutual information} ke representasi terkemas-bit yang menurunkan puncak alokasi CoDA-EDA dari 2.195 KB menjadi 450 KB tanpa mengubah pohon dependensi yang dihasilkan maupun waktu komputasinya.